\documentclass[UTF8,a4paper,12pt]{ctexart}

\usepackage{amsmath,amssymb,bm}
\usepackage{booktabs}
\usepackage{geometry}
\usepackage{longtable}
\usepackage{array}
\usepackage{enumitem}
\usepackage{microtype}
\usepackage{hyperref}

\geometry{left=2.6cm,right=2.6cm,top=2.5cm,bottom=2.5cm}
\setlength{\parindent}{2em}
\setlength{\parskip}{0.35em}
\renewcommand{\arraystretch}{1.35}
\setlist{nosep}
\hypersetup{hidelinks}

\title{方法与实验部分}
\author{}
\date{}

\begin{document}

\maketitle

\setcounter{section}{2}

\section{提出的模型}

本节建立面向大型多线换乘站整体应急疏散的网络模型、队列预测模型和评价模型。模型的基本思路是：先将站内站台、站厅、换乘连接、楼扶梯、闸机和出口组织为可计算的疏散网络，再把多来源客流在共享设施处的到达、排队和服务过程显式纳入路径代价，最后通过统一的容量执行层检验路径方案能否在物理约束下完成疏散。

\subsection{站内疏散网络与共享设施服务容量模型}

\subsubsection{站内空间网络}

将车站公共区抽象为有向图
\[
G=(V,E),
\]
其中，\(V\) 为节点集合，表示线路站台、列车门区、站厅、换乘连接节点、楼梯、扶梯、闸机、出口前区域和出入口；\(E\) 为有向边集合，表示相邻空间或设施之间允许乘客移动的连接关系。边 \(e=(i,j)\in E\) 具有长度 \(l_e\)、宽度或等效宽度 \(w_e\)、方向 \(d_e\)、边类型 \(s_e\)、通行能力 \(c_e\) 和所属线路或区域等属性。节点 \(v\in V\) 具有面积 \(A_v\)、当前占用人数 \(N_v(t)\)、节点类型和空间接收上限等属性。

该网络不是单纯的几何最短路网络。对楼梯、扶梯、换乘通道、闸机和出口前连接等具有独立服务能力或空间接收限制的设施，模型将其作为可被多个来源客流共同占用的资源处理。令 \(R\) 为共享设施资源集合，建立边到资源的映射
\[
\rho:E\rightarrow R\cup\{\varnothing\}.
\]
若 \(\rho(e)=r\)，则边 \(e\) 的进入请求需要消耗资源 \(r\) 的服务能力；若 \(\rho(e)=\varnothing\)，则该边仅受空间密度、下游接收和普通边容量约束。映射到同一资源 \(r\) 的多条上游边共享同一服务能力，不能分别为每条边重复计算容量。由此，来自不同线路站台、站厅和换乘方向的客流在同一楼梯、扶梯、换乘通道或闸机前的竞争关系被纳入同一个资源层。

\subsubsection{多来源客流与换乘来源}

设来源组集合为 \(\mathcal{G}\)。每个来源组 \(g\in\mathcal{G}\) 对应一类具有相同初始位置或到达逻辑的客流，包括各线路站台等待客流、列车到达客流、站厅客流和换乘来源客流。换乘来源客流不作为独立于全站疏散之外的研究对象，而是大型换乘站中影响设施负荷分布的重要来源。其作用体现在两个方面：一是换乘客流的起点和行进方向跨越不同线路公共区；二是换乘客流常与同线路出站客流共同使用楼扶梯、通道、闸机和出口前区域。

乘客需求以自然到达批次表示。一个批次记为
\[
b=(g,n_b,x_b,t_b,P_b,q_b),
\]
其中，\(g\) 为来源组，\(n_b\) 为批次人数，\(x_b\) 为批次当前节点，\(t_b\) 为批次进入路径决策的时间，\(P_b\) 为已确定或候选路径，\(q_b\) 为批次在资源服务队列中的状态。采用批次而非单个行人作为路径决策单元，可以保留列车到达、站台释放和换乘流到达的时间结构，同时避免将数万名乘客逐一进行重复路径搜索。

\subsubsection{设施服务能力与整数容量}

设共享资源 \(r\) 的服务率为 \(\mu_r\)，单位为人每秒。仿真步长为 \(\Delta t\)，上一时间步未形成完整整数人的容量余量为 \(\xi_r(t)\)，则时间步 \(t\) 内资源 \(r\) 可释放的整数服务能力为
\[
C_r(t)=\left\lfloor \mu_r\Delta t+\xi_r(t)\right\rfloor .
\]
释放容量后，余量更新为
\[
\xi_r(t+\Delta t)=\mu_r\Delta t+\xi_r(t)-C_r(t).
\]
当多个上游节点在同一时间步请求进入同一资源 \(r\) 时，模型先在资源层汇总需求，再依据 \(C_r(t)\) 统一分配通过人数。未获得容量的人数保留在上游节点、资源入口队列或空间阻塞状态中，不被计为已经通过该设施。

设施服务率由几何尺寸或设施数量折算。楼梯和通道的容量系数采用公开标准条文中常用的最大通过能力取值：1 m 宽楼梯下行 4200 人/h、上行 3700 人/h、双向混行 3200 人/h；1 m 宽通道单向 5000 人/h、双向混行 4000 人/h。模型实现中运行扶梯和闸机采用的容量值与公开标准表中的推荐值不完全相同，因此在本文中标记为模型采用值，不作为龙阳路站官方实测值。具体参数来源在案例研究中统一列出。

\subsubsection{下游接收与空间回溢}

共享设施能够释放服务能力并不等价于下游空间一定能够接收相应人数。设节点 \(v\) 的空间接收上限为 \(\bar N_v\)，当前占用人数为 \(N_v(t)\)，本时间步已被其他移动请求预留进入的人数为 \(B_v(t)\)，当前请求进入人数为 \(y_v(t)\)，则目标节点接收约束为
\[
N_v(t)+B_v(t)+y_v(t)\leq \bar N_v .
\]
若该约束不满足，即使资源入口仍有剩余服务能力，请求也不能被全部接受。这一处理用于描述高负荷疏散中常见的空间回溢现象：瓶颈不只发生在设施入口，也可能因下游通道、站厅或出口前区域接收不足而向上游扩散。

\subsubsection{速度--密度关系与物理通行时间}

边内行走速度由客流密度决定。本文采用分段速度--密度关系
\[
v(k)=
\begin{cases}
v_f, & k\leq k_f,\\
\max(v_f-ak,0), & k_f<k<k_j,\\
0, & k\geq k_j,
\end{cases}
\]
其中，\(k\) 为边内密度，\(v_f\) 为自由流速度，\(k_f\) 为自由流密度阈值，\(a\) 为速度下降斜率，\(k_j\) 为拥挤密度上限。边 \(e\) 在预计进入时刻 \(\eta\) 的物理行走时间为
\[
\tau_e(\eta)=\frac{l_e}{v(k_e(\eta))}.
\]

当前实现使用 \(v_f=1.427\ \mathrm{m/s}\)、\(k_f=0.2\ \mathrm{p/m^2}\)、\(a=0.3549\ \mathrm{m^2/(p\cdot s)}\)、\(k_j=4.0\ \mathrm{p/m^2}\)，并采用 \(3.0\ \mathrm{p/m^2}\) 作为高密度判定阈值。这组参数来自对“基于改进 A* 算法的多层邮轮疏散系统仿真”中速度--密度模型和 Improved A* 对照方法的复现，本文只将其作为路径代价和密度判定的文献复现参数，不写成龙阳路站实测速度。楼梯和扶梯速度上限属于车站疏散模型中的设施类型设定，待获得现场测量或运营资料后再区分实测值与模型值。

\subsection{设施到达负荷预测与时间依赖路径组织模型}

在多线换乘站内，路径决策面对的关键状态不是“此刻某个设施是否拥挤”，而是“该批客流到达该设施时需要等待多久”。为此，本文将共享设施视为具有服务率和到达事件的时间依赖资源，在路径搜索过程中估计候选批次到达各设施时的队列和空间接收状态。

\subsubsection{确定性到达事件}

对一个候选批次 \(b\)，若其候选路径经过资源 \(r\)，则该批次对 \(r\) 的影响发生在预计到达时刻 \(\eta_{b,r}\)，而不是路径搜索发生的当前时刻。模型维护资源 \(r\) 的确定性到达事件集合
\[
\mathcal{A}_r(t)=\{(\eta_j,n_j)\mid \eta_j\geq t,\ j=1,2,\ldots\},
\]
其中 \(\eta_j\) 为已确定批次进入资源 \(r\) 的预计时刻，\(n_j\) 为对应人数。该集合只记录尚未消耗资源入口容量的整数批次；已在执行层获得资源容量的在途人员不再重复加入同一资源的未来队列。

同一时间步内，批次按照 \(t_b\)、来源组、当前节点和批次编号的确定性顺序处理。某一批次路径确定后，其预计到达资源的整数人数立即写入相应 \(\mathcal{A}_r(t)\)，后续批次的路径搜索即可看到本轮此前已经承诺的设施占用。这一机制使路径组织从同步“看当前状态”转化为按批次滚动更新的“看未来到达状态”。

\subsubsection{设施队列预测}

设资源 \(r\) 在时刻 \(t\) 的实际等待队列为 \(Q_r(t)\)，候选批次预计在 \(\eta\) 时刻到达该资源。定义队列演化算子
\[
\widehat Q_r(\eta)=
\mathcal{Q}\left(Q_r(t),\mathcal{A}_r(t),\mu_r,t,\eta\right),
\]
其中 \(\mathcal{Q}(\cdot)\) 在 \(t\) 到 \(\eta\) 的时间区间内按服务率释放队列，并在各到达事件发生时加入相应批次人数。若 \(\mu_r>0\)，候选批次到达资源 \(r\) 时因前序需求产生的预计等待时间为
\[
W_r(\eta)=\frac{\widehat Q_r(\eta)}{\mu_r}.
\]
对没有共享服务属性的普通移动边，\(W_r(\eta)\) 取 0；对闸机、楼扶梯、换乘通道等共享资源，该项反映批次预计到达时将面对的排队压力。该预测项与当前时刻的可见队列不同，它把路径前段行走时间、前序设施等待和其他批次的已承诺到达共同传递到后续设施。

\subsubsection{时间依赖路径代价}

对批次 \(b\) 的候选路径 \(P=(e_1,e_2,\ldots,e_m)\)，设第 \(i\) 条边的预计进入时刻为 \(\eta_i\)，资源标识为 \(r_i=\rho(e_i)\)。路径总代价定义为
\[
J(P\mid b)=\sum_{i=1}^{m}
\left[\tau_{e_i}(\eta_i)+W_{r_i}(\eta_i)+S_{r_i}(n_b)+W^{\mathrm{sp}}_{e_i}(\eta_i)\right],
\]
其中，\(\tau_{e_i}\) 为密度约束下的物理行走时间，\(W_{r_i}\) 为候选批次到达资源时的预计队列等待，\(S_{r_i}(n_b)\) 为批次自身占用服务资源形成的平均服务时间，\(W^{\mathrm{sp}}_{e_i}\) 为下游空间接收不足造成的预计等待。当前实现采用
\[
S_{r_i}(n_b)=\frac{n_b}{2\mu_{r_i}}
\]
作为批次自身服务时间项；当边不对应有限服务资源或 \(\mu_{r_i}\) 不适用时，该项取 0。

预计进入时刻按路径顺序递推：
\[
\eta_{i+1}=\eta_i+\tau_{e_i}(\eta_i)+W_{r_i}(\eta_i)+S_{r_i}(n_b)+W^{\mathrm{sp}}_{e_i}(\eta_i).
\]
因此，后续设施的等待时间取决于前序路段和设施造成的到达时刻变化。路径评价由此成为时间依赖问题，而不是静态最短路径问题。

\subsubsection{动态重评估与防摆动}

动态路径组织只作用于仍处于明确选择阶段的批次。已经进入在途执行队列的乘客继续完成已经承诺的边，不在设施内部或边内瞬时改路。对可重评估批次，设保留当前路径的剩余成本为 \(C_{\mathrm{stay}}\)，候选切换路径成本为 \(C_{\mathrm{switch}}\)，相对收益为
\[
R=\frac{C_{\mathrm{stay}}-C_{\mathrm{switch}}}{C_{\mathrm{stay}}}.
\]
仅当候选路径有效、下一跳发生变化、批次仍处于允许决策阶段、\(C_{\mathrm{switch}}<C_{\mathrm{stay}}\)，且 \(R>g_{\min}\) 时，模型接受路径切换。当前 Mode 4 场景采用 \(g_{\min}=0.20\)，该值为本研究的防摆动阈值设定，不表述为外部实测或通用校准参数。

\subsection{车站疏散效果综合评价方法}

评价体系围绕全站疏散效率、共享设施瓶颈、来源组差异和微观仿真一致性展开。本文不把某一个单项指标预设为方法有效性的唯一证据，而是在同一容量执行条件下同时考察尾部清空、等待暴露、关键设施负荷和路径分配变化。

\subsubsection{整体疏散效率指标}

设 \(T_p\) 表示累计疏散人数达到总人数 \(p\%\) 的时刻。本文报告 \(T_{95}\)、\(T_{99}\) 和 \(T_{100}\)。其中，\(T_{100}\) 表示最后一名目标乘客完成疏散的时间；\(T_{95}\) 和 \(T_{99}\) 用于刻画尾部客流进入最后清空阶段前的系统效率。

累计停滞人秒定义为
\[
I_{\mathrm{stat}}=\int_0^{T_{100}}N_{\mathrm{stat}}(t)\,\mathrm{d}t,
\]
其中，\(N_{\mathrm{stat}}(t)\) 为时刻 \(t\) 处于等待、阻塞或未能获得移动容量状态的乘客数。进一步报告人均停滞时间、平均移动时间、平均总疏散时间、总移动距离和有效疏散速度。总移动距离和平均移动时间用于识别路径组织是否通过增加绕行换取更少等待。

\subsubsection{设施瓶颈与负荷均衡指标}

对每个共享资源 \(r\) 记录通过人数 \(F_r\)、峰值队列 \(Q^{\max}_r\)、队列人秒 \(I_{q,r}\)、首次排队时刻、最后排队时刻、平均等待和利用率。关键设施拥堵不只由通过人数判断，而由 \(Q^{\max}_r\)、\(I_{q,r}\) 和最后排队时刻共同判断。

设施负荷分布采用 Jain 指数
\[
J(x)=\frac{(\sum_{i=1}^{n}x_i)^2}{n\sum_{i=1}^{n}x_i^2}.
\]
该指标越接近 1，表示负荷分配越均匀。本文分别计算出口负荷 Jain 指数和关键设施负荷 Jain 指数。两类均衡指标只作为解释路径分配变化的诊断量；若出口利用更均衡但关键设施等待增加，不能据此单独认定疏散效果改善。

\subsubsection{来源组与换乘客流指标}

按来源组统计清空时间、出口选择、关键设施通过量和等待暴露。换乘客流作为来源组进入全站评价体系，其结果用于解释不同线路和换乘方向在共享设施上的汇聚机制。本文关注的是换乘客流如何改变全站设施负荷结构，而不是仅提高某一类换乘客流的单独疏散速度。

\section{动态路径组织模型的求解方法}

本节说明如何将上述时间依赖路径代价和共享容量约束转化为可执行的计算流程。两种路径方法在同一疏散网络和同一物理执行层中运行，其差异限定在路径评价和未来队列信息的使用方式。

\subsection{时间依赖路径搜索的可执行化}

\subsubsection{搜索标签}

对批次 \(b\)，AA 的搜索入口为时间依赖 A*。搜索标签写为
\[
L=(v,\eta,g,P,\pi),
\]
其中，\(v\) 为当前节点，\(\eta\) 为预计到达该节点的时刻，\(g\) 为从批次当前位置到 \(v\) 的累计时间代价，\(P\) 为已形成的局部路径，\(\pi\) 为前驱信息。标签从 \(u\) 沿边 \(e=(u,v)\) 扩展时，先计算物理行走时间，再按照预计进入时刻查询资源队列和下游空间接收状态，得到
\[
\eta'=\eta+\tau_e(\eta)+W_{\rho(e)}(\eta)+S_{\rho(e)}(n_b)+W^{\mathrm{sp}}_e(\eta).
\]
静态自由流时间只作为启发式下界，用于提高搜索效率；完整路径代价仍按时间依赖资源状态逐段更新。

\subsubsection{标签保留与剪枝}

若两个标签到达同一节点，其中一个标签在累计到达时刻、累计代价和可行性方面均不优于另一个标签，则删除劣势标签。若两个标签虽然到达同一节点，但对应的预计时刻不同，导致后续共享资源队列状态不同，则保留相应标签。该设置避免把具有不同未来到达状态的路径错误合并为同一个静态节点。

搜索过程中剔除明显回环路径、方向不允许路径和不可达目标。对共享设施容量和下游空间接收的最终约束不在搜索阶段直接兑现，而在执行层统一检验。这样做的原因是：同一时间步内多个来源批次可能同时请求同一资源，只有在全局请求汇总后才能确定实际获得容量的人数。

\subsection{对照方法与本文方法}

\subsubsection{ImprovedAStar 对照方法}

ImprovedAStar 在本文中仅作为计算对照方法。它复现改进 A* 类疏散路径搜索的基本思想，在同一车站网络中依据距离、速度--密度关系和当前拥堵状态计算路径。该方法不使用 AA 的未来资源到达事件索引，不估计候选批次到达设施时的排队状态，也不执行基于相对收益阈值的动态重评估。ImprovedAStar 不代表龙阳路站运营方制定或发布的疏散预案，不能写成“原有疏散方案”。

\subsubsection{AdaptiveQueueAwareAStar}

AdaptiveQueueAwareAStar（AA）以批次为路径决策单元，在搜索过程中将物理行走时间、预计资源等待、批次服务时间和空间接收等待共同纳入路径代价。每个自然到达批次独立执行完整图搜索；前一批次确定路径后，其预计进入共享资源的整数人数立即写入未来到达事件集合，后一批次据此更新候选路径的预计等待。AA 不使用固定 \(K\) 条候选路径缓存，也不把同节点批次直接复用为同一条动态路径。

两种方法共用同一网络、同一客流输入、同一设施服务率、同一速度--密度关系、同一下游接收约束和同一终止判据。由此，结果差异主要来自路径评价是否利用设施到达时刻的队列预测，而不是来自物理执行规则改变。

\subsection{共享容量执行流程}

每个仿真时间步按以下顺序推进：
\begin{enumerate}
    \item 更新已完成边内移动的在途人员，释放其进入目标节点或资源队列；
    \item 根据列车到达、站台释放和换乘来源规则生成自然到达批次；
    \item 对可决策批次调用对应路径方法，得到候选路径和预计资源到达事件；
    \item 将本轮已确定批次的整数到达事件写入资源事件集合；
    \item 汇总所有上游节点对同一共享资源的移动请求，计算 \(C_r(t)\) 并执行资源容量分配；
    \item 检查目标节点接收上限、边内密度和空间回溢，将未被接受的人数保留在上游或队列状态；
    \item 记录疏散人数、来源组位置、资源队列、设施通过量、路径变更和阻塞事件；
    \item 对仍处于允许决策阶段的批次进入下一时间步重评估。
\end{enumerate}

共享容量执行层的作用是保证路径方案和设施服务过程使用同一容量口径。路径搜索只能产生候选方向和预计到达事件，不能直接让乘客越过楼梯、扶梯、闸机或出口。该设置使模型能够区分“路径在网络上可达”和“乘客在容量约束下实际通过”两个层次。

\section{案例研究}

本文以龙阳路站为案例开展模型验证和方法对比。案例研究按照“站点与参数核定--设施到达负荷预测分析--模型方案性能分析--综合疏散效果验证”的顺序展开。

\subsection{案例站点与参数来源}

\subsubsection{站点背景}

截至 2026 年 6 月 30 日，中国内地累计有 58 个城市投运城市轨道交通线路 386 条，运营线路总长度为 13268.30 km；其中上海运营线路长度为 1041.96 km。龙阳路站位于上海轨道交通网络中。上海市交通委员会公开信息显示，18 号线一期北段开通后，龙阳路站可实现 2 号线、7 号线、16 号线、18 号线之间换乘，并可与磁浮线转乘；同一资料还指出该站建筑结构复杂、换乘通道多，预计进出站及换乘客流较大。上述公开资料用于说明案例站点的网络背景和五线换乘属性。

本文模型包含 L2、L7、L16、L18 和 Maglev 五类线路公共区，空间对象包括线路站台、站厅、换乘通道、楼梯、扶梯、闸机、出口前连接区和地面出口。站内几何尺寸、楼扶梯宽度、闸机数量、出口宽度和连接关系以 CAD 图纸与本研究模型配置为基础；在 CAD 图纸完成逐项核对前，正文中统一标注为“CAD/模型配置值”，不写成官方发布数据。

\subsubsection{参数来源分类}

\begin{table}[htbp]
\centering
\caption{案例模型参数来源及正文表述方式}
\small
\begin{tabular}{p{0.18\textwidth}p{0.27\textwidth}p{0.24\textwidth}p{0.21\textwidth}}
\toprule
参数类别 & 主要内容 & 当前来源 & 正文表述方式 \\
\midrule
公开官方资料 & 全国城轨规模、上海线路规模、龙阳路五线换乘关系 & 中国城市轨道交通协会 2026 年上半年快报；上海市交通委员会公开新闻 & 可作为案例背景事实引用 \\
公开标准取值 & 楼梯、通道最大通过能力；站台事故疏散时间要求；最小宽度要求 & 公开标准条文及政府政策解读中列出的地铁设计规范相关条款 & 可作为容量折算依据，但需注明不是车站实测值 \\
CAD/模型配置值 & 站台面积、楼扶梯宽度、通道宽度、闸机数量、出口宽度、节点连接关系 & 龙阳路站 CAD 图纸与本研究代码配置 & 写为 CAD 量测值或模型配置值，待图纸核对后定稿 \\
文献复现参数 & 速度--密度关系、ImprovedAStar 对照参数、高密度阈值 & 蒙盾等 2022 年改进 A* 疏散仿真文献及本研究复现代码 & 写为文献复现/模型判定参数 \\
算法设定参数 & \(\Delta t=1\ \mathrm{s}\)、\(g_{\min}=0.20\)、队列事件登记规则 & 本研究代码与实验配置 & 写为算法设定，并通过敏感性或边界讨论说明影响 \\
情景输入 & Mode 1 的 2187 人、Mode 4 的 17905 人及列车叠加规则 & 本研究场景生成脚本 & 写为情景输入，不写成实测客流 \\
\bottomrule
\end{tabular}
\end{table}

按照上述分类，本文在参数表中区分“实测/CAD量测值”“模型设定值”和“情景假设值”。公开资料中未能核实的设施尺寸不补写为官方数据；公开标准与代码取值不一致的项目，以本研究模型采用值为准，并说明其来源边界。

\subsection{设施到达负荷预测分析}

本节对应模型链条中的中间模块验证，重点检验 AA 是否在路径搜索时使用了批次到达设施时的预计服务状态。分析对象包括楼梯、扶梯、换乘通道、闸机和出口前连接等关键共享资源。

对每个关键资源 \(r\)，记录实际队列 \(Q_r(t)\)、未来确定到达事件 \(\mathcal{A}_r(t)\)、预测队列 \(\widehat Q_r(\eta)\)、预测等待 \(W_r(\eta)\)、实际通过人数和未获容量人数。对同一来源组的代表性批次，分解候选路径代价中的物理行走时间、资源等待、批次服务时间和空间接收等待，比较 ImprovedAStar 与 AA 在相同当前网络状态下对设施负荷的判断差异。

该分析不以“搜索次数增加”作为模型有效性的证据，而是关注三个可验证关系：第一，当前批次的路径评价是否随预计到达时刻改变；第二，前序批次的路径承诺是否进入后续批次的资源队列预测；第三，路径改变是否能够在执行层表现为关键设施到达时段、队列人秒或峰值队列的变化。

\begin{table}[htbp]
\centering
\caption{设施到达负荷预测结果记录表}
\small
\begin{tabular}{p{0.16\textwidth}p{0.16\textwidth}p{0.16\textwidth}p{0.16\textwidth}p{0.16\textwidth}}
\toprule
场景 & 资源类别 & 预测队列误差或趋势 & 最大预测等待 & 对路径选择的影响 \\
\midrule
Mode 1 & 关键楼扶梯/闸机 & 待冻结 & 待冻结 & 待冻结 \\
Mode 4 & 关键楼扶梯/闸机 & 待冻结 & 待冻结 & 待冻结 \\
\bottomrule
\end{tabular}
\end{table}

\subsection{模型方案性能分析}

\subsubsection{实验分组}

本文设置低负荷和高负荷两类疏散场景。Mode 1 为低负荷场景，包含各线路站台等待、站厅和换乘基础客流，场景输入人数为 2187 人。该场景用于检验在设施竞争较弱时，时间依赖队列预测是否引入不必要的绕行、等待或出口偏置。

Mode 4 为高负荷场景，在基础站内客流上叠加双向满载列车客流，场景输入人数为 17905 人，并启用 L2 列车来源按区域划分的客流表示。该场景用于检验在多线路客流集中释放、换乘客流与站台/列车客流共同占用设施的条件下，设施到达负荷预测对尾部清空和关键设施瓶颈的影响。

两种负荷条件均比较 ImprovedAStar 与 AA。二者使用同一网络、同一客流、同一时间步、同一容量折算、同一速度--密度关系、同一空间接收约束和同一停止规则。ImprovedAStar 是算法对照组，AA 是本文提出的时间依赖队列感知路径组织方法。

\subsubsection{整体疏散性能}

整体疏散性能首先报告 \(T_{95}\)、\(T_{99}\)、\(T_{100}\)、累计停滞人秒、人均停滞时间、平均移动时间、平均总疏散时间、有效疏散速度、总移动距离和运行时间。低负荷结果用于判断模型在一般应急客流下是否保持稳定，高负荷结果用于判断模型在多源客流强汇聚下是否改变尾部清空和等待暴露。

\begin{table}[htbp]
\centering
\caption{两种负荷场景下的整体疏散性能比较}
\small
\begin{tabular}{llrrrrrrr}
\toprule
场景 & 方法 & 人数 & \(T_{95}\) & \(T_{99}\) & \(T_{100}\) & 停滞人秒 & 人均停滞 & 总移动距离 \\
\midrule
Mode 1 & ImprovedAStar & 2187 & 待冻结 & 待冻结 & 待冻结 & 待冻结 & 待冻结 & 待冻结 \\
Mode 1 & AA & 2187 & 待冻结 & 待冻结 & 待冻结 & 待冻结 & 待冻结 & 待冻结 \\
Mode 4 & ImprovedAStar & 17905 & 待冻结 & 待冻结 & 待冻结 & 待冻结 & 待冻结 & 待冻结 \\
Mode 4 & AA & 17905 & 待冻结 & 待冻结 & 待冻结 & 待冻结 & 待冻结 & 待冻结 \\
\bottomrule
\end{tabular}
\end{table}

结果解释按指标层级展开：先比较是否完成全部目标乘客疏散，再比较 \(T_{95}\)、\(T_{99}\) 和 \(T_{100}\)，随后解释停滞人秒、移动时间和总移动距离之间的权衡。若 AA 缩短等待但增加行走距离，正文需要同时报告这两个方向的变化，不能只保留有利指标。

\subsubsection{设施瓶颈与来源组分解}

在整体指标之后，进一步报告关键设施队列和来源组分解。对每个场景，按队列人秒和峰值队列识别主要瓶颈设施，比较两种方法下这些设施的到达曲线、服务释放曲线和最后排队时刻。按来源组统计各线路站台、列车来源、站厅来源和换乘来源经过关键设施的人数和等待时间，用于解释尾部疏散差异的站内机制。

\begin{table}[htbp]
\centering
\caption{关键设施负荷与来源组分解结果记录表}
\small
\begin{tabular}{p{0.13\textwidth}p{0.15\textwidth}p{0.16\textwidth}p{0.16\textwidth}p{0.16\textwidth}p{0.13\textwidth}}
\toprule
场景 & 方法 & 主要瓶颈设施 & 队列人秒 & 峰值队列 & 主要来源组 \\
\midrule
Mode 1 & ImprovedAStar & 待冻结 & 待冻结 & 待冻结 & 待冻结 \\
Mode 1 & AA & 待冻结 & 待冻结 & 待冻结 & 待冻结 \\
Mode 4 & ImprovedAStar & 待冻结 & 待冻结 & 待冻结 & 待冻结 \\
Mode 4 & AA & 待冻结 & 待冻结 & 待冻结 & 待冻结 \\
\bottomrule
\end{tabular}
\end{table}

换乘客流分析放在来源组分解中进行。若某一换乘来源在高负荷条件下显著占用特定换乘通道、楼扶梯或闸机，正文将进一步检查该设施是否也是全站尾部清空的限制环节；若换乘来源的路径改变只造成局部排队转移，而没有降低系统停滞或尾部时间，也需要在结果中如实呈现。

\subsection{综合疏散效果验证}

\subsubsection{敏感性分析}

敏感性分析用于检验结论对关键模型参数的依赖程度。当前已纳入的敏感性项为速度--密度模型中的高密度判定阈值，比较约 \(2.0\)、\(3.0\) 和 \(3.9\ \mathrm{p/m^2}\) 等设定下的 \(T_{95}\)、\(T_{100}\)、停滞人秒、关键设施队列和方法间差异。该实验只能说明结论对密度判定阈值的稳健性，不能直接替代对客流规模、设施服务率或到达重叠程度的完整适用边界分析。

\begin{table}[htbp]
\centering
\caption{速度--密度阈值敏感性分析记录表}
\small
\begin{tabular}{p{0.15\textwidth}p{0.16\textwidth}p{0.16\textwidth}p{0.16\textwidth}p{0.16\textwidth}}
\toprule
阈值设定 & 方法 & \(T_{95}\) & \(T_{100}\) & 关键设施队列变化 \\
\midrule
\(2.0\ \mathrm{p/m^2}\) & ImprovedAStar/AA & 待冻结 & 待冻结 & 待冻结 \\
\(3.0\ \mathrm{p/m^2}\) & ImprovedAStar/AA & 待冻结 & 待冻结 & 待冻结 \\
\(3.9\ \mathrm{p/m^2}\) & ImprovedAStar/AA & 待冻结 & 待冻结 & 待冻结 \\
\bottomrule
\end{tabular}
\end{table}

\subsubsection{Pathfinder 微观仿真对照}

Pathfinder 对照用于检查宏观网络模型所得路径组织和瓶颈识别在微观行人仿真中是否具有相近趋势。本文将宏观模型导出的来源组、路径链和路径人数转换为 Pathfinder 中的 Occupant Group 和 Behavior，并在相同几何边界和对应出口设置下运行微观仿真。Pathfinder 官方文档支持行为动作、房间或队列目标、队列服务时间和门选择代价等设置，因此适合用于微观尺度的趋势对照。

Pathfinder 不作为真实观测数据，也不作为宏观模型的绝对真值。对照重点包括：累计疏散曲线形态、尾部清空时段、楼扶梯和闸机附近拥堵区域、出口使用分布、乘客拥堵时间和路径使用差异。若两种尺度结果一致，说明宏观模型识别出的瓶颈和路径组织趋势具有跨尺度一致性；若出现差异，则从导航网格几何、个体避让、门选择规则、队列建模方式和容量映射差异解释。

\begin{table}[htbp]
\centering
\caption{Pathfinder 对照结果记录表}
\small
\begin{tabular}{p{0.16\textwidth}p{0.18\textwidth}p{0.18\textwidth}p{0.18\textwidth}p{0.18\textwidth}}
\toprule
场景 & 方法 & Pathfinder 清空时间 & 拥堵时间/区域 & 与宏观模型趋势关系 \\
\midrule
Mode 1 & ImprovedAStar & 待冻结 & 待冻结 & 待冻结 \\
Mode 1 & AA & 待冻结 & 待冻结 & 待冻结 \\
Mode 4 & ImprovedAStar & 待冻结 & 待冻结 & 待冻结 \\
Mode 4 & AA & 待冻结 & 待冻结 & 待冻结 \\
\bottomrule
\end{tabular}
\end{table}

\subsubsection{综合评价顺序}

案例结果按以下顺序组织：首先报告设施到达负荷预测结果，确认队列预测确实参与路径评价；其次报告 Mode 1 和 Mode 4 的整体疏散指标，说明方法在低负荷和高负荷下的系统表现；再次通过关键设施和来源组分解解释差异来源，重点分析换乘客流与其他来源客流在共享设施上的汇聚关系；随后报告敏感性分析，说明关键参数扰动下结果是否保持；最后用 Pathfinder 微观仿真对宏观趋势进行交叉验证。

\section*{本阶段已核对资料}

\begin{enumerate}
    \item 中国城市轨道交通协会：《2026 年上半年中国内地城轨交通线路概况》，2026 年 7 月 1 日发布。公开页与 PDF 显示，截至 2026 年 6 月 30 日，中国内地累计 58 个城市、386 条运营线路、13268.30 km；上海运营线路长度 1041.96 km。链接：https://www.camet.org.cn/xytj/xxfb/822295470702661.shtml。
    \item 上海市交通委员会：《还有2天！两条轨交线开通初期运营，上海地铁又添“世界第一”》，2021 年 12 月 28 日。资料说明龙阳路站可实现 2、7、16、18 号线换乘并与磁浮线转乘，且车站建筑结构复杂、换乘通道多。链接：https://jtw.sh.gov.cn/xwfb2/20211228/103fc31a81644a72b0c05300f5ecee95.html。
    \item 北京市交通委员会：《北京市城市轨道交通安全运营管理办法》修订条文解读，公开列出地铁设计规范相关条款和车站各部位最大通过能力表，包括楼梯、通道、自动扶梯和检票机能力取值。链接：https://jtw.beijing.gov.cn/xxgk/zcjd/202001/t20200102_1552428.html。
    \item 蒙盾、胡卓、张华军：基于改进 A* 算法的多层邮轮疏散系统仿真，系统仿真学报，2022，34(6):1375--1382，DOI:10.16182/j.issn1004731x.joss.21-0075。本文速度--密度关系和 ImprovedAStar 对照参数按该文献及本研究代码复现。
    \item Thunderhead Engineering Pathfinder User Manual。本文仅将 Pathfinder 作为微观仿真趋势对照工具，相关行为、队列和门选择设置以实际使用版本手册为准。
\end{enumerate}

\end{document}
